\scanpage{77}
% archive foliation: 76
réunion des $\varphi_{S'/S}(\mathcal U)$, où $\mathcal U$ parcourt les
ouverts affines de $\mathfrak X$.

\uline{Corollaire 2} Supposons \deleted{de plus}
\inserted{(que tte partie finie de $\mathfrak X$ soit contenue dans un ouvert affine, et)}
$\mathfrak X$ de type fini \inserted{sur $S'$}, et $S$ loc. noeth.
Alors $\varphi_{S'/S}(\mathfrak X)$ est de type fini sur $S$.

La question étant locale \inserted{sur $S$}, on peut supposer $S$ affine.
Tout résulte du

\uline{Lemme} Soit $S$ noethérien, $\mathfrak X$ de type fini sur $S$,
$n$ un entier. \inserted{\deleted{\unsure{les}} $(\mathfrak X_i)_{i\in I}$ famille d'ouverts},
On suppose que tte partie \deleted{finie} $E$ de $\mathfrak X$ telle que
$\operatorname{card}(E)\leq n$ est contenue dans un ouvert
\deleted{\unsure{affine}} $\mathfrak X_i$. Alors il existe une s.-famille finie
de $(\mathfrak X_i)_{i\in I}$ ayant la $=$ propriété.

En effet, l'hyp. signifie que
\[
\bigcup_{i\in I}\underbrace{\mathfrak X_i\times_S-\times_S\mathfrak X_i}_{n\ \text{facteurs}}
=\underbrace{\mathfrak X\times_S-\times_S\mathfrak X}_{n\ \text{facteurs}}.
\]
On applique le fait que $\mathfrak X\times_S\unsure{-}\times_S\mathfrak X$
est \inserted{un espace} noethérien.

\margin{$\unsure{2|\mathcal Z}$ Le foncteur $\underline W^a$ n'est pas local en tous les
\unsure{cas. Car} \illegible{1} il n'existe pas \illegible{1}}
